\documentclass[lang=cn,11pt]{elegantbook}

\title{工科数学分析}
\subtitle{学习笔记}

\author{夏同}
\institute{华南理工大学}
\date{\today}
\version{1.0}
\bioinfo{自定义}{信息}

\extrainfo{靡不有初，鲜克有终}

\setcounter{tocdepth}{3}

\logo{logo-blue.png}
\cover{751_01.png}

% 本文档命令
\usepackage{array}

\newcommand{\ccr}[1]{\makecell{{\color{#1}\rule{1cm}{1cm}}}}

\begin{document}

\maketitle
\frontmatter

\tableofcontents

\mainmatter

\chapter{微分方程}

\begin{introduction}
  \item 微分方程的定义与阶~\ref{def:ode-def}
  \item 实际问题中的微分方程模型
  \item 微分方程的解~\ref{def:ode-solution}, \ref{def:ode-general-solution}
  \item 定解问题与初值问题~\ref{def:ode-initial-value}
  \item 积分曲线~\ref{def:ode-integral-curve}
  \item 验证与求解方法
\end{introduction}

\section{微分方程的基本概念}

在微积分中，我们研究的函数直接反映了现实世界中量与量之间的关系。然而，面对稍微复杂的实际问题时——例如由切线求曲线方程、物体沿直线的运动、放射性物质的衰变、人口的增长以及加热物体的冷却等——反映运动规律的函数往往不能直接写出，但却比较容易建立这些量及其导数之间的关系式。

这种联系自变量、未知函数以及它们的导数或微分的方程，就称为\textbf{微分方程}（Differential Equation）。微分方程是描述自然界中变化规律的最有力数学工具之一。从天体力学到量子场论，从人口动力学到金融衍生品定价，微分方程几乎渗透到了每一个科学领域。\par

本章将从基本概念出发，逐步建立对常微分方程的系统认识，为后续学习各类具体的解法打下基础。

\begin{definition}[微分方程] \label{def:ode-def}
含有未知函数的导数或微分的方程称为\textbf{微分方程}（Differential Equation）。按照未知函数的个数和类型：
\begin{itemize}
    \item 未知函数是一元函数的方程称为\textbf{常微分方程}(ODE)；
    \item 未知函数是多元函数的方程称为\textbf{偏微分方程}(PDE)。
\end{itemize}
在微分方程中，未知函数的导数或微分的最高阶数称为该方程的\textbf{阶}（Order）。
\end{definition}

\begin{example}[一阶微分方程]
以下均为一阶微分方程：
\[
\frac{\mathrm{d}y}{\mathrm{d}t}=1-y,\quad
\frac{\mathrm{d}y}{\mathrm{d}t}+y=\sin t,\quad
(x^2-y^2)\,\mathrm{d}x+(x^2+y^2)\,\mathrm{d}y=0.
\]
其中前两个为显式方程（可写成 $y' = f(x,y)$ 的形式），第三个为隐式微分方程。
\end{example}

\begin{example}[二阶微分方程]
以下为二阶微分方程：
\[
\frac{\mathrm{d}^2y}{\mathrm{d}t^2}+a^2\sin t=0,\quad
y''+y'^2-y=x.
\]
注意第二个方程因含 $y'^2$ 项而成为非线性方程——线性与非线性之分将在后续章节详细讨论。
\end{example}

\begin{definition}[微分方程的解] \label{def:ode-solution}
若函数 $y=y(x)$ 代入某微分方程后使其变为恒等式，则称此函数为该方程的一个\textbf{解}（Solution）。例如 $y=\mathrm{e}^x$ 是 $\dfrac{\mathrm{d}y}{\mathrm{d}x}=y$ 的解；$s=\cos t$ 和 $s=\sin t$ 均为 $\dfrac{\mathrm{d}^2s}{\mathrm{d}t^2}+s=0$ 的解。
\end{definition}

\begin{definition}[通解] \label{def:ode-general-solution}
若微分方程的解中所含互相独立的任意常数的个数与方程的阶数相同，则称这样的解为该方程的\textbf{通解}（General Solution）。\\
更严格地，对于 $n$ 阶方程，若解 $y=y(x;c_1,\dots,c_n)$ 满足在某邻域内 Wronski 行列式不为零，则 $c_1,\dots,c_n$ 独立且该解为通解。\\
\textbf{反例}：$y=(C_1+2C_2)\sin t$ 中两个常数可合并为一个，不是通解。
\end{definition}

\begin{definition}[特解与定解条件] \label{def:ode-initial-value}
为确定通解中的任意常数而附加的条件称为\textbf{定解条件}。定解条件和微分方程共同构成一个\textbf{定解问题}。满足定解条件的解称为\textbf{特解}。\\
反映系统初始状态的条件称为\textbf{初始条件}。初始条件与微分方程构成的定解问题称为\textbf{初值问题}（Cauchy Problem/IVP），标准形式：\\
一阶：$\begin{cases} y'=f(x,y), \\ y(x_0)=y_0; \end{cases}$~~~~~二阶：$\begin{cases} y''=f(x,y,y'), \\ y(x_0)=y_0,\; y'(x_0)=y_0'. \end{cases}$
\end{definition}

\begin{definition}[积分曲线] \label{def:ode-integral-curve}
微分方程解的图形（即解曲线 $y=y(x)$ 在平面上的图像）称为该方程的\textbf{积分曲线}（Integral Curve）。通解对应于积分曲线族；初值问题的几何意义则是从曲线族中找出经过定点 $(x_0, y_0)$（及指定切线方向）的那条特定曲线。
\end{definition}

\subsection{由通解反求微分方程}

已知含任意常数的函数族（通解），可通过足够多次求导后联立消去常数来反推原方程。

\begin{example}
已知通解 $y=c_1\mathrm{e}^{-x}+c_2\mathrm{e}^{2x}$，求其所满足的微分方程。
\end{example}

\begin{solution}
对 $y$ 求导：
\[
y'=-c_1\mathrm{e}^{-x}+2c_2\mathrm{e}^{2x},~~
y''=c_1\mathrm{e}^{-x}+4c_2\mathrm{e}^{2x}.
\]
$y+y''=5c_2\mathrm{e}^{2x}$ 与 $y'-y=-3c_2\mathrm{e}^{2x}$ 联立消去 $c_2\mathrm{e}^{2x}$：$3(y+y'')+5(y'-y)=0$，整理得
\[
\boxed{y''-y'-2y=0}.
\]
这是二阶常系数线性齐次方程，特征根恰为 $r_1=-1,\; r_2=2$，与出发时的通解结构吻合。
\end{solution}

\begin{note}
\begin{itemize}
\item  通解一定包含所有解吗？不一定。某些方程可能存在\textbf{奇解}（Singular Solution）——不在通解表达式中的孤立解，如 Clairaut 方程 $(y')^2+f(y')=0$ 常有奇解。
\item  所有微分方程都有通解吗？不是。$(y')^2+y^4=0$ 在实范围内仅有零解 $y\equiv 0$，不含任意常数，无传统意义上的通解。此外存在性定理保证了解在一定条件下存在，但未必能用初等函数表示。
\end{itemize}
\end{note}

\section{练习}

\begin{exercise}[相关概念]
回答下列问题：
\begin{itemize}
\item 什么叫微分方程？
\item 微分方程的阶如何定义？
\item 什么叫微分方程的通解？通解中含多少个任意常数？
\item 什么叫初值问题？什么叫微分方程的特解？
\end{itemize}
\end{exercise}
\begin{solution}
\begin{itemize}
    \item 含有未知函数及其导数（或微分）的方程称为微分方程。简而言之，它表达了某个未知函数与其变化率之间的关系。
    \item 微分方程中出现的未知函数的最高阶导数的阶数，称为该微分方程的阶。例如，含有 $\frac{\mathrm{d}^2y}{\mathrm{d}x^2}$ 且没有更高阶导数的方程就是二阶微分方程。
    \item 如果微分方程的解中含有任意常数，且任意常数的个数与微分方程的阶数相同（且这些常数是相互独立的），这样的解称为微分方程的通解。通解中含有的任意常数个数等于微分方程的阶数。
    \item 初值问题：附加了用来确定通解中任意常数的条件（即初始条件）的微分方程求解问题。特解：通过给定的初始条件，确定了通解中的任意常数后所得到的具体的解，不含有任意常数。
\end{itemize}
\end{solution}

\begin{exercise}[验证解]
证明对任意常数 $P_0$，函数 $P(t)=P_0\mathrm{e}^t$ 满足微分方程 $\displaystyle \frac{\mathrm{d}P}{\mathrm{d}t}=P$。
\end{exercise}
\begin{solution}
已知函数为 $P(t)=P_0\mathrm{e}^t$。对该函数关于 $t$ 求导：$$\frac{\mathrm{d}P}{\mathrm{d}t} = \frac{\mathrm{d}}{\mathrm{d}t}(P_0\mathrm{e}^t) = P_0\mathrm{e}^t$$将 $P(t)=P_0\mathrm{e}^t$ 代入上式右端，得到：$$\frac{\mathrm{d}P}{\mathrm{d}t} = P$$左端等于右端，因此函数 $P(t)=P_0\mathrm{e}^t$ 满足该微分方程。证明完毕。
\end{solution}

\begin{exercise}[求参数]
设 $Q=C\mathrm{e}^{kt}$ 满足微分方程 $\displaystyle \frac{\mathrm{d}Q}{\mathrm{d}t}=-0.03Q$，问 $C$ 与 $k$ 为何值？
\end{exercise}
\begin{solution}
对已知函数 $Q=C\mathrm{e}^{kt}$ 关于 $t$ 求导：$$\frac{\mathrm{d}Q}{\mathrm{d}t} = C \cdot k\mathrm{e}^{kt} = k(C\mathrm{e}^{kt}) = kQ$$题目已知该函数满足微分方程：$$\frac{\mathrm{d}Q}{\mathrm{d}t} = -0.03Q$$对比两式可得：$k = -0.03$。由于该方程是一阶线性齐次微分方程，常数 $C$ 作为初始条件决定的系数，在没有给出具体初始条件的情况下，$C$ 可以是任意常数。结论： $k = -0.03$，$C$ 为任意常数。
\end{solution}

\begin{exercise}[求参数]
设 $y=\cos\omega t$，求 $\omega$ 的值，使 $y$ 满足方程 $\displaystyle \frac{\mathrm{d}^2y}{\mathrm{d}t^2}+9y=0$。
\end{exercise}
\begin{solution}
已知 $y=\cos\omega t$。求一阶导数：$$\frac{\mathrm{d}y}{\mathrm{d}t} = -\omega\sin\omega t$$求二阶导数：$$\frac{\mathrm{d}^2y}{\mathrm{d}t^2} = -\omega^2\cos\omega t$$将 $y$ 和二阶导数代入微分方程 $\frac{\mathrm{d}^2y}{\mathrm{d}t^2}+9y=0$ 中：$$-\omega^2\cos\omega t + 9\cos\omega t = 0\implies(9 - \omega^2)\cos\omega t = 0$$为了使上式对任意 $t$ 恒成立，必须满足：$$9 - \omega^2 = 0$$解得：$\omega = \pm 3$。
\end{solution}

\begin{exercise}[求参数]
在下列各题中, 利用所给的初始条件确定函数关系中的常数:
\begin{enumerate}
    \item[(1)] $\displaystyle x^2-y^2=C$, $\displaystyle y|_{x=0}=5$;
    \item[(2)] $\displaystyle y=\mathrm{e}^{2x}(C_1+C_2x)$, $\displaystyle y|_{x=0}=0,\,y'|_{x=0}=1$;
    \item[(3)] $\displaystyle y=C_1\sin(x-C_2)$, $\displaystyle y|_{x=\pi}=1,\,y'|_{x=\pi}=0$.
\end{enumerate}
\end{exercise}
\begin{solution}
\begin{enumerate}
    \item  将 $x=0$ 和 $y=5$ 代入方程：$$0^2 - 5^2 = C \implies C = -25$$方程为：$x^2-y^2=-25$ 或 $y^2-x^2=25$。
    \item 首先利用 $y(0)=0$：$$0 = \mathrm{e}^0(C_1 + C_2 \cdot 0) \implies C_1 = 0$$将 $C_1=0$ 代入原函数，得 $y = C_2 x \mathrm{e}^{2x}$。对其求导：$$y' = C_2(\mathrm{e}^{2x} + 2x\mathrm{e}^{2x})$$利用 $y'(0)=1$：$$1 = C_2(\mathrm{e}^0 + 0) \implies C_2 = 1$$所以常数为 $C_1 = 0$, $C_2 = 1$。（特解为 $y=x\mathrm{e}^{2x}$）
    \item 对原函数求导：$$y' = C_1\cos(x-C_2)$$代入初始条件：$$y(\pi) = C_1\sin(\pi-C_2) = C_1\sin C_2 = 1 \quad \text{(式1)}$$$$y'(\pi) = C_1\cos(\pi-C_2) = -C_1\cos C_2 = 0 \quad \text{(式2)}$$由（式1）知 $C_1 \neq 0$，因此由（式2）必有 $\cos C_2 = 0$。解得 $C_2 = \frac{\pi}{2} + k\pi$ ($k$ 为整数)。为了简便，取 $k=0$，即 $C_2 = \frac{\pi}{2}$。此时 $\sin(\frac{\pi}{2}) = 1$。代入（式1）得 $C_1 \cdot 1 = 1 \implies C_1 = 1$。所以常数可取为 $C_1 = 1$, $C_2 = \frac{\pi}{2}$。（特解为 $y=\sin(x-\frac{\pi}{2}) = -\cos x$）
\end{enumerate}
\end{solution}

\begin{exercise}[函数对应微分方程]
指出下列给出的函数满足哪个微分方程:
\begin{center}
\begin{tabular}{@{}lll@{}}
\toprule
 \textbf{函数} & \textbf{微分方程} \\
\midrule
(1)  $y=2\sin x$ & (a) $\displaystyle \frac{\mathrm{d}y}{\mathrm{d}x}=-2y$ \\[6pt]
(2) $y=\sin2x$ & (b) $\displaystyle \frac{\mathrm{d}y}{\mathrm{d}x}=2y$ \\[6pt]
(3)  $y=\mathrm{e}^{2x}$ & (c) $\displaystyle \frac{\mathrm{d}^2y}{\mathrm{d}x^2}=-y$ \\[6pt]
(4)  $y=\mathrm{e}^{-2x}$ & (d) $\displaystyle \frac{\mathrm{d}^2y}{\mathrm{d}x^2}=-4y$ \\
\bottomrule
\end{tabular}
\end{center}
\end{exercise}
\begin{solution}
    对四个给定函数求导：
    \begin{itemize}
    \item $y=2\sin x \implies y' = 2\cos x \implies y'' = -2\sin x = -y$，故(a) $y' = -2y$ 对应 (4)
    \item $y=\sin2x \implies y' = 2\cos2x \implies y'' = -4\sin2x = -4y$，故(b) $y' = 2y$ 对应 (3)
    \item $y=\mathrm{e}^{2x} \implies y' = 2\mathrm{e}^{2x} = 2y$，故(c) $y'' = -y$ 对应 (1)
    \item $y=\mathrm{e}^{-2x} \implies y' = -2\mathrm{e}^{-2x} = -2y$，故(d) $y'' = -4y$ 对应 (2)
    \end{itemize}
\end{solution}

\begin{exercise}[函数对应微分方程]
指出下列哪个函数是哪个微分方程的解?
\begin{center}
\begin{tabular}{@{}lll@{}}
\toprule
\textbf{函数} & \textbf{微分方程} \\
\midrule
(1)  $y=\mathrm{e}^x$ & (a) $y''-y=0$ \\[6pt]
(2)  $y=x^3$ & (b) $x^2y''+2xy'-2y=0$ \\[6pt]
(3)  $y=\mathrm{e}^{-x}$ & (c) $x^2y''-6y=0$ \\[6pt]
(4)  $y=x^{-2}$ & \\
\bottomrule
\end{tabular}
\end{center}
\end{exercise}
\begin{solution}
逐一验证左侧给出的函数：
\begin{enumerate}
\item $y=\mathrm{e}^x\implies y'=\mathrm{e}^x$, $y''=\mathrm{e}^x$。代入 (a) $y''-y=0$，得 $\mathrm{e}^x-\mathrm{e}^x=0$，成立。对应 (a)。
\item $y=x^3\implies y'=3x^2$, $y''=6x$。代入 (c) $x^2y''-6y=0$，得 $x^2(6x) - 6(x^3) = 6x^3 - 6x^3 = 0$，成立。对应 (c)。
\item $y=\mathrm{e}^{-x}\implies y'=-\mathrm{e}^{-x}$, $y''=\mathrm{e}^{-x}$。代入 (a) $y''-y=0$，得 $\mathrm{e}^{-x}-\mathrm{e}^{-x}=0$，成立。同样对应 (a)。（注：(1)和(3)都是方程(a)的解，构成了它的基础解系）
\item $y=x^{-2}\implies y'=-2x^{-3}$, $y''=6x^{-4}$。代入 (b) $x^2y''+2xy'-2y=0$：$$x^2(6x^{-4}) + 2x(-2x^{-3}) - 2(x^{-2}) = 6x^{-2} - 4x^{-2} - 2x^{-2} = 0$$成立。对应 (b)。
\end{enumerate}
\end{solution}

\begin{exercise}[曲线族对应微分方程]
曲线族常常作为微分方程的解出现. 试将下列曲线与微分方程对应起来:
\begin{center}
\begin{tabular}{@{}lll@{}}
\toprule
 \textbf{曲线族} & \textbf{微分方程} \\
\midrule
(1) $y=x\mathrm{e}^{kx}$ & (a) $\displaystyle \frac{\mathrm{d}y}{\mathrm{d}x}=\frac{y}{x}$ \\[6pt]
(2) $y=x^p$ & (b) $\displaystyle \frac{\mathrm{d}y}{\mathrm{d}x}=\frac{y\ln y}{x}$ \\[6pt]
(3) $y=\mathrm{e}^{kx}$ & (c) $\displaystyle \frac{\mathrm{d}y}{\mathrm{d}x}=\frac{y}{x}\left(1+\ln\frac{y}{x}\right)$ \\[6pt]
(4) $y=mx$ & (d) $\displaystyle \frac{\mathrm{d}y}{\mathrm{d}x}=\frac{y\ln y}{x\ln x}$ \\
\bottomrule
\end{tabular}
\end{center}
\end{exercise}
\begin{solution}
我们需要通过消去曲线族中的参数（$k, p, m$）来得到微分方程：
\begin{enumerate}
    \item[(1)] 对于 $y=x\mathrm{e}^{kx}$，两边对 $x$ 求导得：
    $$ \frac{\mathrm{d}y}{\mathrm{d}x} = \mathrm{e}^{kx} + kx\mathrm{e}^{kx} $$
    由原方程可得 $\displaystyle \mathrm{e}^{kx} = \frac{y}{x}$，两边取对数得 $\displaystyle kx = \ln\frac{y}{x}$。
    将其代入导数表达式中，得：
    $$ \frac{\mathrm{d}y}{\mathrm{d}x} = \frac{y}{x} + \left(\ln\frac{y}{x}\right)\frac{y}{x} = \frac{y}{x}\left(1+\ln\frac{y}{x}\right) $$
    因此，(1) 对应 (c)。
    \item[(2)] 对于 $y=x^p$，两边对 $x$ 求导得：
    $$ \frac{\mathrm{d}y}{\mathrm{d}x} = px^{p-1} = p\frac{x^p}{x} = p\frac{y}{x} $$
    由原方程 $y=x^p$ 两边取自然对数得 $\ln y = p\ln x$，从而得到参数 $\displaystyle p = \frac{\ln y}{\ln x}$。
    将 $p$ 代入导数表达式中，得：
    $$ \frac{\mathrm{d}y}{\mathrm{d}x} = \frac{\ln y}{\ln x} \cdot \frac{y}{x} = \frac{y\ln y}{x\ln x} $$
    因此，(2) 对应 (d)。
    \item[(3)] 对于 $y=\mathrm{e}^{kx}$，两边对 $x$ 求导得：
    $$ \frac{\mathrm{d}y}{\mathrm{d}x} = k\mathrm{e}^{kx} = ky $$
    由原方程两边取自然对数得 $\ln y = kx$，从而 $\displaystyle k = \frac{\ln y}{x}$。
    将 $k$ 代入导数表达式中，得：
    $$ \frac{\mathrm{d}y}{\mathrm{d}x} = \frac{\ln y}{x} \cdot y = \frac{y\ln y}{x} $$
    因此，(3) 对应 (b)。
    \item[(4)] 对于 $y=mx$，两边对 $x$ 求导得：
    $$ \frac{\mathrm{d}y}{\mathrm{d}x} = m $$
    由原方程得 $\displaystyle m = \frac{y}{x}$，直接代入导数表达式得：
    $$ \frac{\mathrm{d}y}{\mathrm{d}x} = \frac{y}{x} $$
    因此，(4) 对应 (a)。
\end{enumerate}
综上所述，对应关系为：(1) --- (c)，(2) --- (d)，(3) --- (b)，(4) --- (a)。
\end{solution}

\begin{exercise}[曲线族对应微分方程]
求抛物线族 $\displaystyle y=C_1(x-C_2)^2$ 的微分方程.
\end{exercise}
\begin{solution}
该曲线族含有两个独立参数 $C_1$ 和 $C_2$，因此我们需要对原式求两次导数，以消去这两个参数，得到一个二阶微分方程。已知：$$y = C_1(x-C_2)^2 \quad \text{(式1)}$$两边对 $x$ 求一阶导数：$$y' = 2C_1(x-C_2) \quad \text{(式2)}$$两边再对 $x$ 求二阶导数：$$y'' = 2C_1 \quad \text{(式3)}$$现在利用这三个式子消去 $C_1$ 和 $C_2$：由（式3）直接得到 $C_1 = \frac{y''}{2}$。将（式2）两边平方，可以巧妙地凑出（式1）的形式：$$(y')^2 = 4C_1^2(x-C_2)^2$$$$(y')^2 = 4C_1 \cdot [C_1(x-C_2)^2]$$观察发现括号内的部分正好是 $y$，于是：$$(y')^2 = 4C_1 y$$最后，把从（式3）得到的 $2C_1 = y''$ 代入上式（即 $4C_1 = 2y''$）：$$(y')^2 = 2y'' y\implies 2yy'' - (y')^2 = 0$$这就是该抛物线族满足的微分方程。
\end{solution}

\section{变量可分离方程及齐次方程}

\begin{introduction}
  \item 变量可分离方程的定义与解法~\ref{def:separable-ode}
  \item 齐次方程及其标准解法~\ref{def:homogeneous-ode}
  \item 可化为齐次方程的类型~\ref{ex:ode-quasi-homog},~\ref{ex:ode-intersect-line}
  \item $R$--$C$ 电路应用~\ref{ex:ode-rc-circuit}
  \item 探照灯反射镜面~\ref{ex:ode-searchlight-mirror}
  \item 马尔萨斯增长模型与逻辑斯蒂模型~\ref{def:malthus-model},~\ref{def:logistic-model}
  \item 牛顿冷却定律与法医应用~\ref{ex:ode-forensic-cooling}
\end{introduction}

接下来我们介绍一阶微分方程的初等解法，其核心思想是把微分方程的求解问题转化为积分问题。需要指出的是，这里并不要求积分一定能用初等函数表示。例如，方程
$\displaystyle \frac{\mathrm{d}y}{\mathrm{d}x}=\frac{y}{x}\left(\mathrm{e}^{\frac{y}{x}}+1\right)$ 的解可以写为
\[
\ln|x| = \mathrm{Li}\!\left(\mathrm{e}^{-\frac{y}{x}}\right) + C,
\]
其中 $\mathrm{Li}(t)=\displaystyle\int\frac{\mathrm{d}t}{\ln t}$ 称为对数积分函数，属于超越函数类。

我们知道，解 $n$ 次代数方程时总希望能用四则运算和方根把解表示出来——但后来证明了对五次及五次以上的代数方程，这种表示一般是不可能的。类似地，对微分方程的求解问题也有同样的限制：\textbf{一般的一阶微分方程没有初等解法}。

尽管如此，若干具有特殊结构的方程类型确实可以通过初等方法求解，且这些类型涵盖了实际应用中出现的相当大一部分情形。本章将系统介绍其中的两类基本形式——变量可分离方程与齐次方程。

\subsection{变量可分离方程}

\begin{definition}[变量可分离方程] \label{def:separable-ode}
形如
\[
\frac{\mathrm{d}y}{\mathrm{d}x}=f(x)\,\phi(y)
\]
的一阶方程称为\textbf{变量可分离方程}，其中 $f(x)$ 和 $\phi(y)$ 分别是 $x$ 和 $y$
的连续函数。
\end{definition}

若 $\phi(y)\neq 0$，则上式可改写为 $\displaystyle \frac{\mathrm{d}y}{\phi(y)}=f(x)\,\mathrm{d}x$，
两边积分得
\[
\int\frac{\mathrm{d}y}{\phi(y)}=\int f(x)\,\mathrm{d}x + C, \tag{5.1}
\]
其中 $C$ 为任意常数。(5.1) 式确定了一个隐函数关系 $y=y(x;C)$，
对其两边微分即可验证它满足原方程。称这样的解为\textbf{隐式通解}。

\begin{remark}[丢失的特解]
若存在常数 $y_0$ 使得 $\phi(y_0)=0$，则直接代入验证可知 $y\equiv y_0$ 也是原方程的一个特解。\textbf{这类特解往往在分离变量过程中丢失，且有时不包含在通解表达式中，必须另行补上才能得到全部解}。
\end{remark}

\begin{example} \label{ex:sep-circle}
求 $\displaystyle\frac{\mathrm{d}y}{\mathrm{d}x}=-\frac{x}{y}$ 的通解。
\end{example}
\begin{solution}
将方程变形为 $y\,\mathrm{d}y=-x\,\mathrm{d}x$，两边积分：
\[
\int y\,\mathrm{d}y=-\int x\,\mathrm{d}x
\;\Longrightarrow\;
\frac{1}{2}y^{2}=-\frac{1}{2}x^{2}+C,
\]
即
\[
\boxed{x^{2}+y^{2}=C},
\]
其中 $C>0$ 为任意常数。这是一族以原点为圆心的同心圆。
\end{solution}

\begin{example}
求方程 $x^{2}y\,\mathrm{d}y+\sqrt{1-y^{2}}\,\mathrm{d}x=0$ 的全部解。
\end{example}
\begin{solution}
当 $1-y^{2}\neq 0$ 时，方程分离为
$\displaystyle -\frac{y\,\mathrm{d}y}{\sqrt{1-y^{2}}}=\frac{\mathrm{d}x}{x^{2}}$，
两边积分得 $\sqrt{1-y^{2}}=-\dfrac{1}{x}+C$，即
\[
\sqrt{1-y^{2}}+\frac{1}{x}=C.
\]

此外，易见 $y=\pm 1$ 也满足原方程（此时 $\sqrt{1-y^2}=0$），
且它们不包含在上面的通解中。因此方程的\textbf{全部解}为
\[
\boxed{\sqrt{1-y^{2}}+\frac{1}{x}=C \quad\text{及}\quad y=\pm 1}.
\]
\end{solution}

\begin{example}
求 $\displaystyle\frac{\mathrm{d}y}{\mathrm{d}x}=y^{2}\cos x$ 满足初始条件 $y(0)=1$ 的特解。
\end{example}
\begin{solution}
当 $y\neq 0$ 时分离变量：$\displaystyle\frac{\mathrm{d}y}{y^{2}}=\cos x\,\mathrm{d}x$，
积分得 $-\dfrac{1}{y}=\sin x + C$，另有特解 $y\equiv 0$。

代入 $y(0)=1$：$-1=0+C$，故 $C=-1$。所求特解为
\[
\boxed{y=\frac{1}{1-\sin x}},\qquad x\neq\frac{\pi}{2}+2k\pi.
\]
\end{solution}

\begin{example} \label{ex:sep-linear-homog}
求 $\displaystyle\frac{\mathrm{d}y}{\mathrm{d}x}=p(x)\,y$ 的通解，其中 $p(x)$ 是连续函数。
\end{example}
\begin{solution}
当 $y\neq 0$ 时分离变量并积分：
\[
\int\frac{\mathrm{d}y}{y}=\int p(x)\,\mathrm{d}x
\;\Longrightarrow\;
\ln|y|=\int p(x)\,\mathrm{d}x+\bar{C},
\]
令 $C=\pm\mathrm{e}^{\bar{C}}$ 得
\[
\boxed{y=C\exp\!\left(\int p(x)\,\mathrm{d}x\right)}.
\]
此外 $y\equiv 0$ 显然也是解，且已包含在上述通解中（取 $C=0$）。\hfill$\square$
\end{solution}

\begin{note}
此例的结果非常重要——它是后续一阶线性方程求解公式的基础。
注意通解中的积分 $\int p(x)\,\mathrm{d}x$ 只需取一个原函数即可，常数部分已吸收到 $C$ 中。
\end{note}

\subsection{齐次方程}

\begin{definition}[齐次方程] \label{def:homogeneous-ode}
若一阶方程 $\displaystyle\frac{\mathrm{d}y}{\mathrm{d}x}=f(x,y)$ 右端的函数满足
$f(tx,ty)=f(x,y)$ 对所有 $t\neq 0$ 成立，
即 $f(x,y)$ 可写成仅依赖于比值 $\dfrac{y}{x}$ 的函数 $\varphi\!\left(\dfrac{y}{x}\right)$，
则称该方程为\textbf{齐次方程}（Homogeneous Equation）。

例如 $\displaystyle\frac{\mathrm{d}y}{\mathrm{d}x}=\frac{xy}{x^{2}+y^{2}}$ 是齐次的，因为
\[
\frac{xy}{x^{2}+y^{2}}
=\frac{\dfrac{y}{x}}{1+\left(\dfrac{y}{x}\right)^{2}}
=: \varphi\!\left(\frac{y}{x}\right).
\]
\end{definition}

对于齐次方程，作代换 $u=\dfrac{y}{x}$（即 $y=ux$），由乘积法则
$\displaystyle\frac{\mathrm{d}y}{\mathrm{d}x}=u+x\frac{\mathrm{d}u}{\mathrm{d}x}$，
代入原方程得
\[
u+x\frac{\mathrm{d}u}{\mathrm{d}x}=\varphi(u)
\quad\Longrightarrow\quad
\frac{\mathrm{d}u}{\mathrm{d}x}=\frac{\varphi(u)-u}{x}. \tag{5.2}
\]
这是一个关于新未知函数 $u(x)$ 的\textbf{变量可分离方程}！

\begin{property}[齐次方程的完整解法]
\begin{enumerate}
  \item 若 $\varphi(u)-u\neq 0$，则 (5.2) 分离为
        $\displaystyle\frac{\mathrm{d}u}{\varphi(u)-u}=\frac{\mathrm{d}x}{x}$，
        积分后回代 $u=\dfrac{y}{x}$ 即得通解；
  \item 若 $\varphi(u)-u=0$ 有实根 $u_1,u_2,\dots,u_m$，
        则直线族 $y=u_i\,x$（$i=1,\dots,m$）也是解，一般不包含在通解中；
  \item 若 $\varphi(u)-u\equiv 0$，则 $\dfrac{\mathrm{d}u}{\mathrm{d}x}\equiv 0$，
        解为一族过原点的直线 $y=cx$。
\end{enumerate}
\end{property}

\begin{example}
求 $\displaystyle\frac{\mathrm{d}y}{\mathrm{d}x}=\frac{y}{x}+\tan\frac{y}{x}$ 的通解。
\end{example}
\begin{solution}
令 $u=\dfrac{y}{x}$，代入化简得 $\displaystyle\frac{\mathrm{d}u}{\mathrm{d}x}=\frac{\tan u}{x}$。
分离变量：$\dfrac{\cos u\,\mathrm{d}u}{\sin u}=\dfrac{\mathrm{d}x}{x}$，两边积分：
$\ln|\sin u|=\ln|x|+C_1$，即 $\sin u=Cx$。
此外 $\tan u=0$（即 $\sin u=0$）对应 $C=0$ 的情形，已被包含。
回代 $u=\dfrac{y}{x}$ 得通解
\[
\boxed{\sin\frac{y}{x}=Cx}.
\]
\end{solution}

\begin{example}
求解 $(x+y)\,\mathrm{d}x-(y-x)\,\mathrm{d}y=0$。
\end{example}
\begin{solution}
令 $y=ux$，则 $\mathrm{d}y=x\,\mathrm{d}u+u\,\mathrm{d}x$，代入化简得
$(1+2u-u^2)\,\mathrm{d}x+x(1-u)\,\mathrm{d}u=0$。分离变量并积分：
$x^2(1+2u-u^2)=C$。回代并补充 $u=1\pm\sqrt{2}$ 的特解（取 $C=0$ 已包含）：
\[
\boxed{x^2+2xy-y^2=C}.
\]
\end{solution}

\begin{example}
求解 $x\,\mathrm{d}y-y\,\mathrm{d}x=\sqrt{x^{2}+y^{2}}\,\mathrm{d}x$。
\end{example}
\begin{solution}
设 $x>0$（$x<0$ 推导完全类似），原方程等价于
$\displaystyle\frac{\mathrm{d}y}{\mathrm{d}x}=\frac{y}{x}+\sqrt{1+\left(\frac{y}{x}\right)^{2}}$。

令 $u=\dfrac{y}{x}$，分离变量得 $\displaystyle\frac{\mathrm{d}u}{\sqrt{1+u^2}}=\frac{\mathrm{d}x}{x}$，
积分：$\ln\!\left(u+\sqrt{1+u^2}\right)=\ln x+C_1$，即 $u+\sqrt{1+u^2}=Cx$（$C>0$）。
消去根号后化简可得
\[
\boxed{y=\frac{1}{2}\left(Cx^{2}-\frac{1}{C}\right)}.
\]
\end{solution}

\subsection{可化为齐次方程的类型}

\begin{example} \label{ex:ode-quasi-homog}
求解方程 $\displaystyle\frac{\mathrm{d}y}{\mathrm{d}x}=\frac{a_1x+b_1y+c_1}{a_2x+b_2y+c_2}$（系数均为常数）。
\end{example}
\begin{solution}
分三种情形：
\begin{itemize}
\item $c_1=c_2=0$：已是齐次方程，令 $u=y/x$ 即可。
\item $\begin{vmatrix}a_1&b_1\\a_2&b_2\end{vmatrix}=0$（即 $a_1/a_2=b_1/b_2=k$）：
分子写成 $k(a_2x+b_2y)+c_1$，令 $u=a_2x+b_2y$ 化为可分离方程。
\item 行列式非零且 $c_1,c_2$ 不全为零：
两直线有唯一交点 $(\alpha,\beta)$。作平移 $X=x-\alpha,\;Y=y-\beta$，
方程变为 $\displaystyle\frac{\mathrm{d}Y}{\mathrm{d}X}=g(Y/X)$，即化为齐次方程。
\end{itemize}
\end{solution}

\begin{example} \label{ex:ode-intersect-line}
求解 $\displaystyle\frac{\mathrm{d}y}{\mathrm{d}x}=\frac{x-y+1}{x+y-3}$。
\end{example}
\begin{solution}
联立 $\begin{cases}x-y+1=0\\x+y-3=0\end{cases}$ 得交点 $(1,2)$。
令 $X=x-1,\;Y=y-2$ 化为 $\dfrac{\mathrm{d}Y}{\mathrm{d}X}=\dfrac{X-Y}{X+Y}$，
再令 $Y=uX$ 分离变量积分，最终回代化简得
\[
\boxed{x^{2}-y^{2}+2x+6y-2xy=C}.
\]
\end{solution}

\begin{example}
求解初值问题
$\begin{cases}(x^2+2xy-y^2)\mathrm{d}x+(y^2+2xy-x^2)\mathrm{d}y=0,\\ y(1)=1.\end{cases}$
\end{example}
\begin{solution}
两边同除以 $x^2$ 后显式化为齐次形式。令 $y=ux$，
经有理分式分解 $\dfrac{u^2+2u-1}{(u+1)(u^2+1)}=\dfrac{2u}{u^2+1}-\dfrac{1}{u+1}$
积分化简得通解 $\dfrac{x^2+y^2}{x+y}=C$。
代入 $y(1)=1$ 得 $C=1$，故特解为
\[
\boxed{x^{2}+y^{2}=x+y}.
\]
\end{solution}

\begin{example}
求解：(1) $f(xy)y\,\mathrm{d}x+g(xy)x\,\mathrm{d}y=0$；
\quad (2) $\displaystyle\frac{\mathrm{d}y}{\mathrm{d}x}=\frac{1}{x+y}$。
\end{example}
\begin{solution}
\textbf{(1)} 令 $u=xy$，则 $\mathrm{d}u=x\,\mathrm{d}y+y\,\mathrm{d}x$。代入得
$\ln|x|+\displaystyle\int\frac{g(u)}{u[f(u)-g(u)]}\,\mathrm{d}u=C$。\\
\textbf{(2)} 令 $u=x+y$，则 $\dfrac{\mathrm{d}u}{\mathrm{d}x}-1=\dfrac{1}{u}$，
分离变量积分得
\[
\boxed{y-\ln|x+y+1|=C}\quad\text{或}\quad \boxed{x=C\mathrm{e}^{y}-y-1}.
\]
\end{solution}

\begin{remark}
方程中出现 $f(xy),\,f(x\pm y),\,f(x^2\pm y^2),\,f(y/x)$ 等复合结构时，
通常应尝试相应代换 $u=xy,\,x\pm y,\,x^2\pm y^2,\,y/x$ 等，
将方程化为可分离变量型。
\end{remark}

\subsection{应用实例}

\subsubsection*{$R$--$C$ 电路充电过程}

\begin{example} \label{ex:ode-rc-circuit}
如图所示的 $R$--$C$ 电路，电容 $C$ 初始无电荷。
合上开关 K 后电池 $E$ 对电容充电，求电压 $u_C$ 随时间的变化规律。
\end{example}

\begin{solution}
由基尔霍夫第二定律（KVL）：$u_C + RI = E$。
利用 $Q=Cu_C$ 和 $I=\mathrm{d}Q/\mathrm{d}t=C\,\mathrm{d}u_C/\mathrm{d}t$ 得
$RC\,\dfrac{\mathrm{d}u_C}{\mathrm{d}t}+u_C=E$。

这是可分离变量方程。结合初始条件 $u_C(0)=0$，解得
\[
\boxed{u_C(t)=E\left(1-\mathrm{e}^{-\frac{t}{RC}}\right)}. \tag{5.3}
\]
定义时间常数 $\tau=RC$：当 $t=3\tau$ 时 $u_C\approx 0.95E$，
工程上通常认为充电基本完成。\hfill$\square$
\end{solution}

\subsubsection*{探照灯反射镜面的形状}

\begin{example} \label{ex:ode-searchlight-mirror}
制造探照灯反射镜时要求点光源发出的光经镜面反射后平行射出，
试求反射镜面的几何形状。
\end{example}

\begin{solution}
取光源为坐标原点，$x$ 轴平行于反射方向。设镜面由曲线 $y=f(x)$ 绕 $x$ 轴旋转而成。
根据光学反射定律（入射角等于反射角），导出微分方程
\[
\frac{\mathrm{d}y}{\mathrm{d}x}=\frac{y}{x+\sqrt{x^{2}+y^{2}}}. \tag{5.4}
\]

这是齐次方程。采用代换 $v=x/y$ 将 (5.4) 分离变量并积分（计算略），最终得到
\[
\boxed{y^{2}=C(C+2x)},\quad C>0.
\]

这是一条抛物线！因此反射镜面为旋转抛物面 $y^{2}+z^{2}=C(C+2x)$。
这正是抛物线的经典光学性质——来自焦点的所有光线经反射后均沿对称轴方向平行射出。\hfill$\square$
\end{solution}

\subsection{微分方程建模：增长、衰减与冷却}

本节通过三个典型实际模型展示微分方程描述自然规律的威力。

\subsubsection*{马尔萨斯人口增长模型}

\begin{definition}[马尔萨斯模型 / Malthusian Growth] \label{def:malthus-model}
设某孤立物种（无迁入迁出）在时刻 $t$ 的种群数量为 $P(t)$。
假设净增长率（出生率减死亡率）为不随时间和规模变化的常数 $a$，则
\[
\frac{\mathrm{d}P}{\mathrm{d}t}=aP,\quad P(t_0)=P_0. \tag{5.5}
\]
\end{definition}

由例 \ref{ex:sep-linear-homog} 立即得到解
\[
\boxed{P(t)=P_0\,\mathrm{e}^{a(t-t_0)}}. \tag{5.6}
\]
种群呈\textbf{指数增长}（或衰减，若 $a<0$）。

\begin{note}
马尔萨斯模型（1798）在预测短期、小规模种群时效果良好。历史数据表明世界人口约每 34.6 年翻一番，与模型预测高度吻合。然而长期外推却导致荒谬结论：按此速率到 2670 年地球上的人口密度将使人们互相踩着肩膀站成两层。问题出在模型忽略了\textbf{种内竞争}：个体成员之间会为有限资源而竞争。
\end{note}

\subsubsection*{逻辑斯蒂增长模型（Logistic Model）}

为修正马尔萨斯模型的缺陷，引入与 $P^2$ 成正比的竞争项 $-bP^2$（$b>0$）：

\begin{definition}[逻辑斯蒂模型] \label{def:logistic-model}
\[
\frac{\mathrm{d}P}{\mathrm{d}t}=aP-bP^{2}=aP\left(1-\frac{b}{a}P\right). \tag{5.7}
\]
常数 $a,b$ 称为生命系数（Vital Coefficients）。
\end{definition}

通过部分分式分解积分可得显式解：

\begin{property}[逻辑斯蒂模型解及其极限行为]
满足 $P(t_0)=P_0$ 的解为
\[
\boxed{P(t)=\frac{aP_0}{bP_0+(a-bP_0)\,\mathrm{e}^{-a(t-t_0)}}}. \tag{5.8}
\]
关键特征：
\begin{itemize}
  \item 当 $P_0<a/b$ 时 $P(t)$ 单调递增；当 $P_0>a/b$ 时单调递减；
  \item 无论初值如何，均有 $\displaystyle\lim_{t\to\infty}P(t)=\frac{a}{b}$。
\end{itemize}
极限值 $K:=a/b$ 称为环境的\textbf{承载量}（Carrying Capacity）。
\end{property}

\subsubsection*{牛顿冷却定律与法医应用}

\begin{definition}[牛顿冷却定律] \label{def:newton-cooling}
温度为 $T(t)$ 的物体置于恒温 $m$ 的介质中，变化率正比于温差：
\[
\frac{\mathrm{d}T}{\mathrm{d}t}=-k(T-m),\quad k>0. 
\]
\end{definition}
这个方程的通解前文已给出：$T(t)=m+C\mathrm{e}^{-kt}$。

\begin{example}[法医学中的死亡时间推断] \label{ex:ode-forensic-cooling}
尸体初始体温 $37^\circ\mathrm{C}$。2 小时后降至 $35^\circ\mathrm{C}$，
室温恒定 $20^\circ\mathrm{C}$。下午 4 点发现时体温为 $30^\circ\mathrm{C}$。推断死亡时间。
\end{example}
\begin{solution}
设 $H(t)$ 为尸体温度（$^\circ\mathrm{C}$），$t$ 从死亡时刻起算（小时）。由牛顿冷却定律：
$\dfrac{\mathrm{d}H}{\mathrm{d}t}=-k(H-20),\; H(0)=37$。
分离变量求解得 $H(t)=17\,\mathrm{e}^{-kt}+20$。 利用 $H(2)=35$ 确定 $k$：$35=17\,\mathrm{e}^{-2k}+20$
$\;\Rightarrow\; k\approx 0.063$。
将 $H=30$ 代入 (5.10)：$30=17\,\mathrm{e}^{-0.063t_*}+20$
$\;\Rightarrow\; t_*\approx 8.4$ 小时。
从下午 4 点往前推约 8.4 小时，死亡时间大约在\textbf{上午 7 点 36 分左右}。\hfill$\square$
\end{solution}

\section{练习}




\end{document}